\chapter{Topology and limits}

\section{Topological Spaces}
\label{S:topology}


\begin{definition}
\label{D:topological space}
    A \emph{topological space} is a structure $(X, \mathscr{N}_x)_{x\in X}$  that consists of a set $X$ and, for each point $x\in X$, a collection $\mathscr{N}_x$ of subsets of $X$ that satisfy the following properties:

    \begin{enumerate}
        \item $x\in U$ for every $U\in \mathscr{N}_x$;
        \item If $U,V\in \mathscr{N}_x$, then there exists $W\in \mathscr{N}_x$ such that $W\subseteq U\cap V$;
        \item For every $U\in \mathscr{N}_x$ and every $y\in U$, there exists $V_y\in \mathscr{N}_y$ such that $V_y\subseteq U$.
        
    \end{enumerate}
\end{definition}




Next we introduce an immediate result from the previous definition. The reader should note that Property (3) and the next proposition are equivalent.

\begin{proposition}
\label{P:basic nhoods are open}
    Let $(X, \mathscr{N}_x)_{x\in X}$ be a topological space. For every $U\in \mathscr{N}_x$ and every $y\in U$, there exists $V_y\in \mathscr{N}_y$ such that \[U=\bigcup\limits_{y\in U}V_y\text{.}\]
\end{proposition}

\begin{proof}
  Fix a point $x\in X$, and let $U\in \mathscr{N}_x$ be given.
  By Property (3) above, we have that for every $y\in U$, there is some $V_y\in \mathscr{N}_y$ such that $V_y\subseteq U$.
  Thus, the union of all such $V_y$ is contained in $U$.
  Also, by Property (1), we see that for every $y\in U$, we have $y\in V_y\subseteq \bigcup\limits_{y\in U}V_y$. Thus, $U=\bigcup\limits_{y\in U}V_y$.
\end{proof}

\begin{definition}
  Let $( X, \mathscr{N}_x)_{x\in X}$ be a topological space.

  For each $x\in X$, we refer to $\mathscr{N}_x$ as the \emph{neighborhood base} for $x$. 
  The indexed family $(\mathscr{N}_x)_{x\in X}$ is called a \emph{local neighborhood base} on $X$.
  For all $x\in X$, each set $U\in \mathscr{N}_x$ is called a \emph{basic neighborhood} of $x$. A \emph{neighborhood} of $x$ is any superset of a basic neighborhood of $x$.
    
\end{definition}

\begin{definition}
  A set $O\subseteq X$ is called \emph{open} if for every $x\in O$, there exists $V_x\in \mathscr{N}_x$ such that 
  \[
  O=\bigcup\limits_{x\in O}V_x\text{.}
  \]
  
  A set $C\subseteq X$ is called \emph{closed} if $X\setminus C$ is open. Note that $C$ is closed if for every $x\in X$, and every basic neighborhood of $x$ that intersects $C$, then $x\in C$.
\end{definition}

Note that, by Proposition~\ref{P:basic nhoods are open}, a subset of $X$ is open if and only it is a union of basic neighborhoods.

\begin{proposition}
\label{P:properties of open sets}

    In any topological space $(X, \mathscr{N}_x)_{x\in X}$, the following properties hold:

    \begin{enumerate}
        \item Both $\emptyset$ and $X$ are open subsets of $X$ (and evidently also closed);
        \item Any union of open subsets of $X$ is open;
        \item Any finite intersection of open subsets of $X$ is open.
        
    \end{enumerate}
\end{proposition}

\begin{proof}

\leavevmode

\begin{enumerate}
\item The empty set is open since any condition that must be satisfied by the points in $\emptyset$ is true vacuously. The set $X$ itself is open since, for each point $x\in X$, we may choose $V_x\in \mathscr{N}_x$ arbitrarily so that $X=\bigcup\limits_{x\in X}V_x$.

\item Given any open subsets $(O_i)_{i\in I}$ of $X$, we have that each $O_i$ is a union of basic neighborhoods. Thus, the union $\bigcup\limits_{i\in I}O_i$ is a larger union of basic neighborhoods, and hence also an open subset of $X$.

\item If $O_1$ and $O_2$ are open subsets of $X$, then for each $x\in O_1\cap O_2$, there are $U,V\in \mathscr{N}_x$ such that $U\subseteq O_1$ and $V\subseteq O_2$. By Property (2) of Definition~\ref{D:topological space}, there's $W_x\in \mathscr{N}_x$ such that \[W_x\subseteq U\cap V\subseteq O_1\cap O_2\text{.}\] Thus, $O_1\cap O_2$ is the union of all such $W_x\in \mathscr{N}_x$ that satisfy $W_x\subseteq O_1\cap O_2$ as $x$ ranges over $O_1\cap O_2$.
\end{enumerate}
\end{proof}

Note that an infinite intersection of open sets is not necessarily open, as will be shown in Example \ref{E:basic topological spaces}(a).

\begin{definition}
Let $(X, \mathscr{N}_x)_{x\in X}$ be a topological space.

The \emph{topology} $\mathscr{T}$ on $X$ generated by the local neighborhood base $(\mathscr{N}_x)_{x\in X}$  is the collection of all open subsets of $X$; that is, 
\[
\mathscr{T}= 
\{ O\subseteq X \colon O \text{ is open}\}\text{.}
\]
\end{definition}

Note that two different local neighborhood bases on $X$ can generate the same topology. For instance, examples (a) and (b) below generate the same topology on $\mathbb R$.

\begin{example} Basic examples
\label{E:basic topological spaces}

\begin{enumerate}[(a)]

\item Let $X=\mathbb{R}$, and for each $x\in \mathbb{R}$,  let $\mathscr{N}_x=\{(x-\varepsilon , x+\varepsilon):\varepsilon >0\}$.

We will show that there is an infinite intersection of open subsets of $X$ that is not open. Take, for instance, the infinite intersection \[\bigcap\limits_{n\in \mathbb{N}}(0,1+\frac{1}{n})=(0,1]\text{.}\]
This interval is not open since there is no basic neighborhood $(1-\varepsilon,1+\varepsilon)$ of $1\in X$ that is contained in $(0,1]$.

\item Let $X=\mathbb{R}$, and for each $x\in \mathbb{R}$,  let $\mathscr{N}_x=\{(\alpha , \beta) \colon \alpha<x<\beta, \alpha,\beta\in\mathbb{R} \}$.

\item Let $X=\mathbb{R}^2$, and for each $x\in \mathbb{R}^2$, let $\mathscr{N}_x$ consist of the set of all open disks about $x$ of radius $\varepsilon$ for all $\varepsilon >0$.

\item Let $X$ be any set, and for each $x\in X$, let $\mathscr{N}_x = \{\{x\}\}$; i.e., the only basic neighborhood of $x$ is the singleton $\{x\}$. The resulting topology is called the \emph{discrete} topology on $X$.

Note that, in this topological space, every subset of $X$ is open (since every set is a union of singletons). Hence, every subset is closed as well.

\item Let $X$ be any set, and for each $x\in X$, let $\mathscr{N}_x = \{X\}$. The resulting topology is called the \emph{trivial} (or \emph{indiscrete}) topology on $X$. Note that, in this topological space, only two subsets are open; namely, $\emptyset$ and $X$.

\end{enumerate}

\end{example}

\begin{exercise}
Let $(A,\le)$ be  a linearly ordered set (see Definition~\ref{D:posets} and the exercises that follow it). For $a,b\in A$ with $a<b$, define the \emph{open interval} $(a, b)$ by
\[
(a, b)=\{x\in A \colon a<x<b\}.
\]
Now for $x\in A$, define 
\[
\mathscr{N}_x = \left\{ \,  (a, b) \colon a<x<b \, \right\}.
\]
\begin{enumerate}[(a)]
\item
Prove that $(A, \mathscr{N}_x )_{x\in A}$ is a topological space, i.e., show that $\mathscr{N}_x$ satisfies properties (1)--(3) of the definition of topological space, for all $x\in A$.(Definition~\ref{D:topological space}).
\item
Suppose above that we replace ``linearly ordered'' by ``partially ordered.'' Is $(A, \mathscr{N}_x )_{x\in A}$ a topological space? Why or why not?
\item
Suppose that we define
\[
\mathscr{N}_x = \left\{ \, [x, b) \colon x<b  \, \right\}.
\]
Is $(A, \mathscr{N}_x )_{x\in A}$ a topological space?  Why or why not?
\end{enumerate}
\end{exercise}


\begin{exercises}\hfill


\begin{enumerate}

\item
For any given set $A$, let  $\BbR^A$ denote the set of all functions from $A$ into $\BbR$. We define a topology on $\BbR^A$.  To this effect, fix $f\in\BbR^A$ in order to define a set $\mathscr{N}_f$ of basic neighborhoods of $f$. 

Given $x_1,\dots,x_n \in A$ and $\varepsilon > 0$, define
\[
V_f(x_1,\dots, x_n,\varepsilon)= 
\left\{g\in \BbR^A \colon | g(x_i)-f(x_i) |<\epsilon \text{ for all } i\in \{ 1,\dots, n\} \right\}.
\]
Define now
\[
\mathscr{N}_f = \left\{ V_f(x_1,\dots, x_n,\varepsilon) \colon \varepsilon >0, x_1,\dots,x_n\in A, n\in\mathbb{N} \right\}.
\]
Prove that  $(\BbR^A, \mathscr{N}_f )_{f\in \BbR^A}$ is a topological space, i.e., show that $\mathscr{N}_f$ satisfies properties (1)--(3) of the definition of topological space (Definition~\ref{D:topological space}).


\item 
 Let $(X, \mathscr{N}_x )_{x\in X }$ be a topological space, let $A$ be a set, and denote by  $X^A$ the set of all functions from $A$ into $X$. Define a topology on $X^A$ as follows.  Given $f\in X^A$,    $x_1,\dots,x_n \in A $ and  $V_1\dots, V_n\subseteq X$ such that $V_i\in \mathscr{N}_{f(x_i)}$ for  $i\in \{ 1,\dots, n\}$, define
\[
V_f(x_1,\dots, x_n,V_1,\dots, V_n)= 
\left\{g\in X^A \colon g(x_i)\in V_{f(x_i)}  
\text{ for all } i\in \{ 1,\dots, n\} \right\}.
\]
Then define 
\[
\mathscr{N}_f = \left\{ V_f(x_1,\dots, x_n,V_1,\dots, V_n) \colon \text{$V_i\in \mathscr{N}_{f(x_i)}$ for  $i\in \{ 1,\dots, n\}$} \right\}.
\]
Prove that  $(X^A, \mathscr{N}_f )_{f\in X^A}$ is a topological space, i.e., show that $\mathscr{N}_f$ satisfies properties (1)--(3) of the definition of topological space (Definition~\ref{D:topological space}).

\end{enumerate}

\end{exercises}


\section{Topological Limits}

We now transition to discuss the notion of limit in a topological space. Fix a topological space $(X,\mathscr{N}_x)_{x\in X}$ and an indexed family $(a_i)_{i\in I}$ in $X$. We want to express what it means for $(a_i)_{i\in I}$ to converge to a point $c\in X$. Intuitively, this means that $(a_i)_{i\in I}$ \emph{clusters} around a point $c\in X$; that is, for every basic neighborhood $V_c\in \mathscr{N}_c$, there exists a ``large'' $L\subseteq I$ such that $\{a_i \colon i\in L\}\subseteq V_c$. To clarify what is meant by ``large,'' we proceed by defining the notion of filter, which may be thought of as a collection of all the ``large'' subsets $L$ of $I$.

\begin{definition}
\label{D: filter}

Fix a nonempty set $I$. A \emph{filter} on $I$ is a collection $\mathscr{F}$ of subsets of $I$ that satisfy the following properties:

\begin{enumerate}
        \item $\emptyset\not\in\mathscr{F}$ and $I\in \mathscr{F}$;
        \item If $A\in\mathscr{F}$ and $B\supseteq A$, then $B\in\mathscr{F}$;
        \item If $A,B\in \mathscr{F}$, then $A\cap B\in\mathscr{F}$.
        
    \end{enumerate}
\end{definition}

Similarly, we may define the notion of ``small'' in the following way.

\begin{definition}
\label{D: ideal}

Fix a nonempty set $I$. An \emph{ideal} on $I$ is a collection $\mathscr{J}$ of subsets of $I$ that satisfy the following properties:

\begin{enumerate}
        \item $\emptyset\in\mathscr{J}$ and $I\not\in \mathscr{J}$;
        \item If $A\in\mathscr{J}$ and $B\subseteq A$, then $B\in\mathscr{J}$;
        \item If $A,B\in \mathscr{J}$, then $A\cup B\in\mathscr{J}$.
        
    \end{enumerate}

\end{definition}

Observe that given a filter $\mathscr{F}$ on a nonempty set $I$, we can define an ideal $\mathscr{J}$ on $I$ by \[\mathscr{J}=\{I\setminus A \colon A\in\mathscr{F}\}\text{.}\]
Similarly, given an ideal $\mathscr{J}$ on $I$, we see that \[\mathscr{F} = \{I\setminus A \colon A\in \mathscr{J}\}\] is a filter on $I$.

\begin{notation}
\label{N:Large sets}
If $\mathscr{F}$ is a filter on a nonempty set $I$, then the elements of $\mathscr{F}$ are called \emph{$\mathscr{F}$-large}. A property of elements of $I$ is said to hold for \emph{almost every $i\in I$} or \emph{almost everywhere} in $I \pmod{\mathscr{F}}$ if the property holds for an $\mathscr{F}$-large subset of $I$.
\end{notation}

We are now ready to define the notion of convergence.

\begin{definition}
Let $(X, \mathscr{N}_x)_{x\in X}$ be a topological space, and let $(a_i)_{i\in I}$ be an indexed family of elements in $X$. If $\mathscr{F}$ is a filter on $I$, we say $a_i$ \emph{converges} to a point $c\in X$ with respect to $\mathscr{F}$, denoted $a_i\overset{\mathscr{F}}\to c$, if and only if for every $V_c\in\mathscr{N}_c$, we have $\{i \colon a_i\in V_c\}\in\mathscr{F}$.


We call $c$ an \emph{$\mathscr{F}$-limit} of $(a_i)_{i\in I}$ if and only if $(a_i)_{i\in I}$ converges to $c$ with respect to $\mathscr{F}$.

\end{definition}

In the context of a sequence of real numbers, this is the definition we are familiar with. Consider Example \ref{E:basic topological spaces}(a); namely, the topological space $(\mathbb{R},\mathscr{N}_x)_{x\in \mathbb{R}}$ where $\mathscr{N}_x = \{(x-\varepsilon,x+\varepsilon) \colon \varepsilon>0\}$. Let $I=\mathbb{N}$, and let $\mathscr{F}$ consist of all cofinite subsets of $\mathbb{N}$. Then, if $(a_n)_{n\in \mathbb{N}}$ is a sequence of real numbers, we see that $a_n \overset{\mathscr{F}}\to c$ if and only if for every $V_c\in\mathscr{N}_c$, we have $a_n\in V_c$ for all but finitely many $n\in\mathbb{N}$; that is, for all $\varepsilon >0$, there exists  $N\in\mathbb{N}$ such that for all $n>N$, we have $a_n\in(c-\varepsilon,c+\varepsilon)$.

\begin{definition}
The filter that consists of all cofinite subsets of any nonempty set $I$ is called the \emph{Fr\'echet filter on $I$}.
\end{definition}

Recall that, in the topological space $(\mathbb{R},\mathscr{N}_x)_{x\in \mathbb{R}}$ mentioned above, limits are unique. However, this is not true in every topological space. Consider the indiscrete topology discussed in Example~\ref{E:basic topological spaces}(e); namely, given any set $X$, the indiscrete topology on $X$ is the set $\mathscr{T}=\{\emptyset, X\}$. Given any indexed family $(a_i)_{i\in I}$ in $X$ and any filter $\mathscr{F}$ on $I$, since $X$ is the only basic neighborhood of any point in $X$ and \[\{i\colon a_i\in X\}=I\in\mathscr{F}\text{,}\] it follows that every point in $X$ is a limit of $(a_i)_{i\in I}$.

Thus, a new question arises: What property must a topological space have to ensure that every limit is unique?
This question leads us to the following definition.

\begin{definition}
A topological space $(X,\mathscr{N}_x)_{x\in X}$ is called \emph{Hausdorff} if any two distinct points in $X$ can be separated by disjoint neighborhoods; i.e., if $x,y\in X$ such that $x\not =y$, then there exist basic neighborhoods $V_x\in\mathscr{N}_x$ and $V_y\in\mathscr{N}_y$ such that $V_x\cap V_y=\emptyset$.
\end{definition}

\begin{proposition}
\label{P:Hausdorff->unique limits}
Let $(X,\mathscr{N}_x)_{x\in X}$ be a Hausdorff topological space, let $(a_i)_{i\in I}$ be an indexed family in $X$, and let $\mathscr{F}$ be a filter on $I$.
If $a_i\overset{\mathscr{F}}\to c$ and $a_i\overset{\mathscr{F}}\to d$, then $c=d$.
\end{proposition}

\begin{proof}
Suppose, on the contrary, that 
\begin{align*}
a_i\overset{\mathscr{F}}\to c \text{,}\\
a_i\overset{\mathscr{F}}\to d\text{,}
\end{align*}
and $c\not=d$. Since $(X,\mathscr{N}_x)_{x\in X}$ is Hausdorff, we can separate $c$ and $d$ by disjoint neighborhoods, say, $U$ and $V$. Then we have
\begin{align*}
A&=\{i \colon a_i\in U\}\in\mathscr{F} \text{ and}\\
B&=\{i \colon a_i\in V\}\in\mathscr{F}\text{,}
\end{align*}
so by Property (3) of Definition~\ref{D: filter}, we have $A\cap B=\emptyset\in \mathscr{F}$, which contradicts Property (1). We conclude $c=d$.
\end{proof}


The converse of Proposition~\ref{P:Hausdorff->unique limits} is also true. Although it is not needed for the material in these notes, it can be proved easily, so we have included a proof in Proposition~\ref{P:unique limits->Hausdorff}

\begin{notation}
In case $\mathscr{F}$-limits are unique in a topological space $(X,\mathscr{N}_x)_{x\in X}$ and an indexed family $(a_i)_{i\in I}$ in $X$ converges to $c\in X$, we write 
\[
c=\mathscr{F}\lim a_i \quad \text{or} \quad c= \lim\limits_{i,\mathscr{F}}a_i \text{.}
\]
\end{notation}


\subsection{Principal Ultrafilters \nopunct}


\begin{definition}
An \emph{ultrafilter} $\mathscr{U}$ on a nonempty set $I$ is a filter on $I$ such that for every $A\subseteq I$, either $A\in \mathscr{U}$ or $I\setminus A\in \mathscr{U}$.
\end{definition}

Filters can also be seen as generalized measures. An ultrafilter $\mathscr{U}$ on a set $I$ divides the subsets of $I$ into two categories: $\mathscr{U}$-large and $\mathscr{U}$-small --- the large subsets of $I$ are those in $\mathscr{U}$, and the small subsets of $I$ are those whose complement is in $\mathscr{U}$. Thus, ultrafilters are ``complete'' measures in the sense that every set is measurable. Filters, in contrast, need not be complete in this sense: For an arbitrary filter $\mathscr{F}$ on a set $I$, there may be subsets of $I$ that are neither $\mathscr{F}$-large nor $\mathscr{F}$-small (for instance, for the Fr\'{e}chet filter on $\mathbb N$, any subset of $\mathbb N$ that is neither finite nor cofinite --- e.g., the set of even numbers --- is neither $\mathscr{F}$-large nor $\mathscr{F}$-small.)


\begin{example} Principal Ultrafilters

Fix a nonempty set $I$ and an element $a\in I$. Define 
\[
\mathscr{U}_a=\{A\subseteq I \colon a\in A\}.
\]
Note that $\mathscr{U}_a$ is an ultrafilter on the set $I$. Any ultrafilter defined in this way is called a \emph{principal} (or ``trivial'' or ``fixed'') ultrafilter. More precisely, $\mathscr{U}_a$ is called the principal ultrafilter on $I$ \emph{generated by $a$}.
\end{example}

The preceding example suggests that ultrafilters represent an abstract notion of ``point." Principal ultrafilters are ``hard" points, while nonprincipal ultrafilters are ``soft" points. 

\begin{exercise}
Prove that if $I$ is finite, then every ultrafiter on $I$ is principal.
\end{exercise}




\subsection{Nonprincipal Ultrafilters \nopunct}


\begin{proposition}
An ultrafilter $\mathscr{U}$ on a set $I$ is principal if and only if $\mathscr{U}$ contains a finite subset of $I$.
\end{proposition}

\begin{proof}
If $\mathscr{U}$ is principal, then $\mathscr{U}$ contains a finite set; namely, the singleton that generates it. Suppose, conversely, that $\mathscr{U}$ is nonprincipal. We prove, by induction on $n$, that $\mathscr{U}$ cannot contain $n$-element subsets of $I$ for $n=1, 2,\dots$. The case $n=1$ is given by the definition of nonprincipal.  Suppose that $\mathscr{U}$ contains a set $\{a_1,\dots, a_n\}$, where $a_1,\dots, a_n$ are distinct elements of $I$. Since, by the non nonprincipality assumption, we cannot have $\{a_1\}\in \mathscr{U}$, we must have $I\setminus \{a_1\}\in \mathscr{U}$. But then, since $\mathscr{U}$ is closed under finite intersections, $\mathscr{U}$ must contain $(I\setminus \{a_1\})\cap\{a_1,\dots, a_n\}=\{a_2,\dots, a_n\}$, which contradicts the induction assumption.
\end{proof}

The preceding proposition can be re-stated as follows:

\begin{proposition}
\label{P:nonprincipality and Frechet}
An ultrafilter $\mathscr{U}$ on an infinite set $I$ is nonprincipal if and only if it extends the Fr\'echet filter on $I$.
\end{proposition}

In light of the preceding proposition, to prove that nonprincipal ultrafilters exist, we only have to prove that the Fr\'echet filter on $I$ can be extended to an ultrafilter. The following propositions show that, indeed,  any collection of subsets of $I$ that has the finite intersection property can be extended to an ultrafilter; see, in particular, Corollaries~\ref{C:ultrafilter theorem} and~\ref{C:nonprincipal ultrafilters exist}.

\begin{definition}
\label{D:fip}
A collection $\mathscr{B}$ of subsets of $I$ has the \emph{finite intersection property} (or is \emph{finitely consistent}) if whenever $A_1,\dots, A_n\in\mathscr{B}$, we have 
\[
A_1\cap\dots\cap A_n\neq \emptyset \text{.}
\]
\end{definition}


The finite intersection property is ubiquitous in mathematics. Mathematicians often refer to it as ``f.i.p."

\begin{proposition}
\label{P:filter generated by a base}
Every collection of subsets of $I$ that has the finite intersection property can be extended to a filter on $I$.
\end{proposition}

\begin{proof}
Let $\mathscr{B}$ be a collection of subsets of $I$ with the finite intersection property. Without loss of generality, we can assume that $\mathscr{B}$ is closed under finite intersections. But then, once we have assumed this, to extend $\mathscr{B}$ to a filter, we only have to close it under supersets; in other words, the collection
\[
\mathscr{F}=\{A\subseteq I \colon A\supseteq B 
\text{ for some $B\in \mathscr{B}$} \}
\]
is clearly a filter on $I$.
\end{proof}

Due to the preceding proposition, any collection $\mathscr{B}$ of subsets of $I$ with the finite intersection property is often called a \emph{base} for a filter (or a \emph{filter base}) on $I$. The filter $\mathscr{F}$ constructed in the proof of the Proposition~\ref{P:filter generated by a base} is called the filter \emph{generated} by $\mathscr{B}$.

Now we want to improve Proposition~\ref{P:filter generated by a base} by showing that every filter can be extended to an ultrafilter on $I$.

\begin{proposition}
\label{P:ultrafilter lemma}
Let $\mathscr{F}$ be a filter on a set $I$. Then,  for every $A\subseteq I$, there exists a filter $\mathscr{F}'\supseteq \mathscr{F}$ on $I$ such that either $A\in \mathscr{F}'$ or $I\setminus A \in \mathscr{F}'$
\end{proposition}

\begin{proof}
Fix a set $A\subseteq I$. By Proposition~\ref{P:filter generated by a base}, it suffices to verify that either $\mathscr{F}\cup\{A\}$ or $\mathscr{F}\cup\{I\setminus A\}$ has the finite intersection property. But this is immediate because if $\mathscr{F}\cup\{A\}$ fails to have the finite intersection property, it is because there exist $A_1,\dots, A_n\in \mathscr{F}$ such that
$A_1 \cap\dots \cap A_n\cap A=\emptyset$, and this means that 
$I\setminus A\supseteq A_1 \cap\dots \cap A_n$,
which implies that $\mathscr{F}\cup\{I\setminus A\}$ has the finite intersection property.
\end{proof}


\begin{corollary}[{The Ultrafilter Theorem [A. Tarski, 1930]}]
\label{C:ultrafilter theorem}
Every filter on a set $I$ can be extended to an ultrafilter on $I$.
\end{corollary}

\begin{proof}
Fix a filter $\SF$ on $I$,  let $\alpha=|\SP(I)|$ and let $(A_\xi)_{\xi<\alpha}$ be a list of all the subsets of $I$. By induction, we construct an increasing chain 
\[
\SF=\SF_0\subseteq \SF_1\subseteq \dots\subseteq \SF_\xi\subseteq \dots\qquad (\xi\le\alpha)
\]
of filters in $I$  such that:
\begin{enumerate}
\item If $\xi = \zeta +1$, then 
\[
A_\zeta\in \SF_\xi
\quad\text{or}\quad
I\setminus A_\zeta\in \SF_\xi;
\]
\item If $\xi\le\alpha$ is a nonzero limit ordinal, then
\[
\SF_\xi = \bigcup_{\zeta <\xi} \SF_\zeta\text{.}
\]
\end{enumerate}

The construction is as follows. Fix $\xi<\alpha$ and that $\SF_\zeta$ has been defined for all $\zeta<\xi$. 

If $\xi$ is a successor ordinal, say $\xi=\zeta+1$, by Proposition~\ref{P:ultrafilter lemma},
either $\SF_\zeta \cup \{A_\zeta\}$ or $\SF_\zeta \cup \{I\setminus A_\zeta\}$ has the finite intersection property (it is possible that both sets have the f.i.p.).
We pick one of these sets that has the finite intersection property, and let $\SF_\xi$ be a filter extending it. 

If $\xi$ is a limit ordinal, we simply let $\SF_\xi = \bigcup_{\zeta<\xi}\SF_\zeta$. Since $(\SF_\zeta)_{\zeta<\xi}$ is an increasing chain of filters, it is easy to see that $\SF_\xi$ is closed under finite intersections and supersets (for if $A,B\in \SF_\xi$, then $A$ and $B$ are in some member of the chain), and does not contain the empty set, i.e., $\SF_\xi$ is a filter. 

By construction, $\SF_\alpha$ is an ultrafilter containing $\SF_0=\SF$.

\end{proof}

\begin{corollary}
\label{C:nonprincipal ultrafilters exist}
Nonprincipal ultrafilters exist on any infinite set.
\end{corollary}

\begin{proof}
Let $I$ be an infinite set. By Proposition \ref{P:ultrafilter lemma}, the Fr\'echet filter on $I$ can be extended to an ultrafilter. By Proposition~\ref{P:nonprincipality and Frechet}, such an ultrafilter must be nonprincipal. 
\end{proof}

\begin{corollary}
\label{C:base extended to ultrafilter}
Let $\SB$ be a family of subsets of $I$ with the finite intersection property. Then there exists an ultrafilter $\SU$ on $I$ such that every set in $\SB$ is $\SU$-large.
\end{corollary}

\begin{proof}
By Proposition~\ref{P:filter generated by a base} and Corollary~\ref{C:ultrafilter theorem}.
\end{proof}











